% UC-071: Prototype (practical part); installation and manual are in the appendix (UC-061).
% TODO(verify): add a verified design-science reference to docs/literature.md before citing
% one here; the framing below is stated without citation.
\section{Prototype}
\label{sec:prototype}

The prototype is the artefact of this design-oriented study: it makes the concept of
Section~\ref{sec:concept} usable by a commuter and serves as the object of the evaluation.
Its design follows three requirements derived from the problem. First, the answer must be
about the passenger's own trip, not about the network. Second, every statement must be
checkable against the data, so the interface shows how true each sentence is and how much
data it rests on. Third, the assistant must not invent anything, so no language model is
called at run time.

\subsection{Architecture}
\label{sec:architecture}
The prototype is a web application built with Streamlit on top of the same library that
produces the evaluation results. The data pipeline ingests the daily Ist-Daten files,
validates them, builds a table of measured train arrivals and reconstructs the connections
at the three study hubs (Zürich HB, Bern and Olten) from pairs of arriving and departing
trains. Hub stops are matched by their official identifiers rather than by display names, so
the crisp and the fuzzy view always use the same arrivals. The assistant library contains
the vocabulary, the truth engine, the quality measures, the summary miner, the decision
engine and the crisp official view. The user interface only renders what the library
returns: a separate, tested rendering layer derives every text and number on the screen from
the answer object, so the interface cannot show a number the engine did not compute.

\subsection{User Interface}
\label{sec:ui}
The commuter enters the question in a form: the hub, the arriving and the departing line, the
time of day, the days and their time to change, chosen with a slider. Station names and time
bands come directly from the study configuration, so that the questions match the data
exactly. The answer is organised in three views (Appendix~\ref{app:manual} and
Figures~\ref{fig:app-form}--\ref{fig:app-chatbot}).

\subsubsection*{Answer} A coloured badge shows the decision word. Below it, under \emph{Why},
each explanation sentence is shown with a bar for its truth, coloured by how true the
sentence is (green, amber or red), the number of observations it is based on and its
coverage. The colour of the bar describes the sentence, not the risk: a fully true sentence
that explains a risky answer has a green bar next to a red badge.

\subsubsection*{Official view} The passenger's verdict and its sentence are shown side by side
with the official punctuality verdict and its crisp sentences. This view is the
citizen-facing form of SRQ2: when the two verdicts differ, the passenger sees that their own
time to change is what makes the difference.

\subsubsection*{Chatbot view} For questions of the evaluation set, the view shows the answer
that the language model wrote, stored from the offline runs, with a coloured mark per claim
from the fuzzy checker and a grey mark for claims that cannot be checked. This view is the
citizen-facing form of SRQ3.

Every view ends with the data source and the month of the measurements. The assistant is not
branded as a product of any operator, because the data covers all operators.

\subsection{Design Choices}
\label{sec:design}
Several choices follow from the requirement that every statement be checkable.
\begin{itemize}
  \item \emph{One sentence about the passenger, always.} Every answer contains a sentence
  about the passenger's own connections, either with the best-fitting quantifier or, in the
  vocabulary gap, with the measured share. A risky verdict without any statement about the
  passenger's connections would ask the passenger to trust the decision blindly.
  \item \emph{Truth next to the sentence, not hidden.} Showing the truth and the number of
  observations lets a passenger see that ``few of your connections are missed'' rests on many
  trains or on a handful.
  \item \emph{Both views on one screen.} The official verdict is not hidden but shown next to
  the passenger's verdict, with its own sentences, so the passenger can see why the two
  differ.
  \item \emph{No generation at run time.} The chatbot view shows stored answers from the
  evaluation only. A citizen-facing tool that generates new sentences would need the
  checker in the loop for every sentence, which is outside the scope of this prototype.
\end{itemize}

\subsection{Example}
\label{sec:example}
On the synthetic test data shown in the figures, a commuter at Zürich HB arriving with the
IC1 in the morning peak and three minutes to change receives the answer \emph{risky}: about
half of the observed connections would have been missed. With ten minutes to change, the
same trains give the answer \emph{safe}. The official view, which ignores the time to
change, rates the same arrivals \emph{risky}, because about half of them are three or more
minutes late. The figures will be replaced by screenshots of the real data once the full
evaluation has run.

\subsection{Verification in the Prototype}
\label{sec:proto-verification}
Each answer the prototype shows can be verified with the same function that the evaluation
uses as a correctness gate (Section~\ref{sec:explanation}). The rendering layer is tested to
show only numbers that appear in the answer, and the application is tested end to end on the
synthetic data. Installation from a clean machine takes five commands
(Appendix~\ref{app:install}).
